% ============================================================
%  Literature Review — Task Offloading in IoT–Fog–Cloud Systems
%  Author: Imed Eddine Benmadi
%  Deadline: April 15, 2026
% ============================================================

\documentclass[12pt,a4paper]{article}

\usepackage[a4paper, margin=2.5cm, headheight=15pt]{geometry}
\usepackage{amsmath,amssymb}
\usepackage{graphicx}
\usepackage{booktabs}
\usepackage{array}
\usepackage{tabularx}
\usepackage{hyperref}
\usepackage[numbers,sort&compress]{natbib}
\usepackage{microtype}
\usepackage{parskip}
\usepackage{titlesec}
\usepackage{fancyhdr}
\usepackage{enumitem}
\usepackage{caption}
\usepackage{float}

\hypersetup{
  colorlinks=false,
  hidelinks,
  pdftitle={Task Offloading in IoT-Fog-Cloud: A Literature Review},
}

\pagestyle{fancy}
\fancyhf{}
\rhead{\small Task Offloading in IoT--Fog--Cloud: A Literature Review}
\lhead{\small \thepage}
\renewcommand{\headrulewidth}{0.4pt}

\titleformat{\section}{\large\bfseries}{\thesection}{1em}{}[\vspace{0.1em}\rule{\linewidth}{0.4pt}]
\titleformat{\subsection}{\normalsize\bfseries}{\thesubsection}{1em}{}
\titleformat{\subsubsection}{\normalsize\bfseries\itshape}{\thesubsubsection}{1em}{}

\newcolumntype{L}[1]{>{\raggedright\arraybackslash}p{#1}}
\newcolumntype{C}[1]{>{\centering\arraybackslash}p{#1}}

% ─────────────────────────────────────────────────────────────────────────────
\begin{document}

% ── TITLE PAGE ───────────────────────────────────────────────────────────────
\begin{titlepage}
  \centering
  \vspace*{2cm}
  {\small\textsc{Literature Review --- Distributed Systems Lab}}\\[1.0cm]
  {\LARGE\bfseries
    Optimized Task Offloading in\\[0.4em]
    IoT--Fog--Cloud Systems:\\[0.5em]
    A Comprehensive Literature Review of\\[0.4em]
    Metaheuristic Optimization, Deep\\[0.3em]
    Reinforcement Learning, SDN, and\\[0.3em]
    Mobility-Aware Computation Offloading
  }\\[1.4cm]
  \rule{0.55\linewidth}{0.5pt}\\[1.2cm]
  {\normalsize
    \textbf{Author:} Imed Eddine Benmadi\\[0.5cm]
    \textbf{Keywords:} Task offloading, IoT, Fog computing, NSGA-II,
    deep reinforcement learning,\\
    SDN, DAG task modeling, vehicular networks,
    mobility-aware offloading, proactive handoff
  }\\[1.4cm]
  \rule{0.55\linewidth}{0.5pt}\\[1.0cm]
  {\normalsize April 2026}
  \vfill
  {\small\textit{Prepared for: Distributed Systems Lab --- Literature Review Assignment}}
\end{titlepage}

% ── ABSTRACT ─────────────────────────────────────────────────────────────────
\begin{abstract}
\noindent
The rapid growth of Internet of Things (IoT) deployments in smart-city and
vehicular environments has intensified the demand for low-latency, energy-efficient
computation at the network edge. Task offloading, the process of delegating
compute-intensive work from resource-constrained IoT devices to fog or cloud
infrastructure, has become the central mechanism for meeting these demands
within the three-tier IoT--Fog--Cloud hierarchy.

This paper surveys ten representative works published between 2021 and 2025,
spanning five research directions: multi-objective metaheuristic optimization
using NSGA-II with Minimax Differential Evolution; task pre-filtering via the
TOF-Broker Execution Cost framework; software-defined networking for programmable
traffic management; microservice-oriented DAG task modeling; and mobility-aware
offloading with deep reinforcement learning and proactive handoff prediction.
For each direction, we identify the state of the art, its key contributions, and
what it leaves unresolved.

From this analysis we extract five open research gaps. The two primary gaps
concern integration and cold-start: no existing system simultaneously handles
multi-objective optimization, task pre-filtering, DAG structure, online DRL
adaptation, SDN routing, and mobility; and no system has validated offline
evolutionary pre-training as a cold-start remedy in a full vehicular IoT context.
We argue that bridging these gaps requires a hybrid architecture combining
offline Pareto optimization with online reinforcement learning, applied to
DAG workloads under GPS-driven mobility constraints.
\end{abstract}

\vspace{0.5cm}\hrule\vspace{0.7cm}
\tableofcontents
\newpage

% =============================================================================
\section{Introduction}
\label{sec:intro}
% =============================================================================

The Internet of Things ecosystem is growing faster than most projections anticipated.
Tens of billions of connected devices are expected to be active by 2030~\cite{jafari2021},
and a large share of them operate under conditions that traditional computing
architectures were not designed for: they move, they generate data continuously,
and the decisions they need support with are time-critical. A smart vehicle that
needs to know whether a pedestrian is crossing its path cannot wait 200\,ms for a
response from a distant cloud server. An industrial robot operating on a
millisecond control loop cannot tolerate network jitter. The latency introduced
by conventional cloud computing, typically 50--200\,ms due to wide-area propagation,
makes it structurally unsuitable for these use cases~\cite{jafari2021}.

Fog computing was introduced to close this gap by placing intermediate compute
nodes one or two hops from the end device, bringing latency down to 1--10\,ms.
Together with the cloud, this creates the three-tier hierarchy that now defines
the field: IoT devices at the bottom, fog nodes in the middle, cloud at the top.
The fundamental question this hierarchy raises is \emph{task offloading}: for
each piece of computation arriving from an IoT device, which tier should execute
it, under what constraints, and at what energy cost? This problem is formally
NP-hard~\cite{mashal2024drl}, and at the scale of real deployments no exact
algorithm can solve it in time to be useful.

Five distinct challenges make task offloading hard in practice:

\begin{enumerate}[leftmargin=*, label=\textbf{C\arabic*.}]
  \item \textbf{Boulder problem.} When compute-heavy tasks are routed to fog
    nodes without filtering, they consume the node's resources disproportionately
    and trigger queue cascade failures~\cite{nouri2024tof}.

  \item \textbf{Cold-start problem.} Deep reinforcement learning agents begin
    with random weights and produce near-random decisions for thousands of steps
    before converging, which is unacceptable in deployed systems~\cite{vfc_drl2024}.

  \item \textbf{Static optimization problem.} Evolutionary algorithms like
    NSGA-II take hundreds of generations to converge, making them too slow for
    millisecond real-time decisions~\cite{zhu2025dag}.

  \item \textbf{Network congestion problem.} Reactive SDN routing cannot
    respond to traffic bursts fast enough to prevent packet loss~\cite{gasmi2024sdn,csfd2025}.

  \item \textbf{Mobility problem.} In vehicular networks, devices cross fog
    coverage boundaries faster than tasks complete, causing mid-computation
    handoff failures~\cite{paper5_mobility,mavto2024}.
\end{enumerate}

The central observation of this review is that no single published system resolves
all five challenges together. Individual works advance one or two of them, sometimes
at the expense of the others. The research field is mature enough that each
challenge has been addressed in isolation; what is missing is a unified framework.

This paper is structured as follows. Section~\ref{sec:background} establishes
the formal definitions used throughout. Section~\ref{sec:evolution} traces
how the field has developed across four generations of work. Section~\ref{sec:sota}
surveys the current state of the art in each sub-domain. Section~\ref{sec:gaps}
defines five open research gaps. Section~\ref{sec:proposed} sketches the
architecture that follows from closing them. Section~\ref{sec:discussion}
compares this architecture against related work. Section~\ref{sec:conclusion}
concludes.

% =============================================================================
\section{Background and Formal Definitions}
\label{sec:background}
% =============================================================================

\subsection{Three-Tier IoT--Fog--Cloud Architecture}

The three-tier model organizes computation into three layers. \textbf{Layer~1}
(IoT devices) consists of resource-constrained task generators: smart vehicles,
sensors, drones. They have limited CPU, battery, and memory. \textbf{Layer~2}
(fog nodes) are edge servers co-located with 5G access points, offering
low-latency compute within a coverage radius of roughly 150--300\,m. \textbf{Layer~3}
(cloud) provides high-throughput batch computation at the cost of 20--50\,ms
additional WAN delay.

\subsection{Task Offloading Formulation}

Let $\mathcal{T} = \{t_1, \ldots, t_N\}$ be a batch of $N$ tasks from IoT devices.
Each task $t_i$ is described by:
\begin{equation}
  t_i = \langle \ell_i,\; d_i,\; D_i,\; \mathcal{E}_i \rangle
\end{equation}
where $\ell_i$ is computational demand in Million Instructions (MI), $d_i$ is
input data size in kilobytes, $D_i$ is the end-to-end deadline in milliseconds,
and $\mathcal{E}_i$ is device transmission energy in joules. The offloading
decision assigns each task to a destination
$x_i \in \{\mathrm{FOG}_1, \ldots, \mathrm{FOG}_K, \mathrm{CLOUD}\}$.
The two optimization objectives are:
\begin{align}
  f_1(\mathbf{x}) &= \sum_{i=1}^{N} E_i(x_i) \quad\text{(total energy, minimize)}\\
  f_2(\mathbf{x}) &= \sum_{i=1}^{N} L_i(x_i) \quad\text{(total latency, minimize)}
\end{align}
subject to deadline constraints $L_i(x_i) \leq D_i$ and capacity constraints
$\sum_{i:\,x_i=k} \ell_i \leq C_k$, where $C_k$ is the MIPS capacity of fog node~$k$.

\subsection{Execution Cost and Task Classification}

The Execution Cost (EC) metric introduced in~\cite{nouri2024tof} measures whether
a task is feasible on a fog node:
\begin{equation}
  \mathrm{EC}(t_i) = \frac{\ell_i}{\mathrm{fog\_MIPS}} \quad\text{(seconds)}
  \label{eq:ec}
\end{equation}
A task is a \emph{pebble} if $\mathrm{EC}(t_i) < \theta$ and a \emph{boulder}
if $\mathrm{EC}(t_i) \geq \theta$, where $\theta = 1.0$\,s is the standard
threshold. Boulders occupy a fog node for over a second at full capacity, which
is long enough to starve hundreds of concurrent lightweight tasks.

\subsection{Pareto Dominance}

Solution $\mathbf{x}^{(a)}$ dominates $\mathbf{x}^{(b)}$, written
$\mathbf{x}^{(a)} \prec \mathbf{x}^{(b)}$, when:
\begin{equation}
  f_k(\mathbf{x}^{(a)}) \leq f_k(\mathbf{x}^{(b)})\;\forall k,
  \quad\text{and}\quad
  f_j(\mathbf{x}^{(a)}) < f_j(\mathbf{x}^{(b)})\;\text{for some }j
\end{equation}
The Pareto front $\mathcal{P}^*$ contains all non-dominated solutions. No element
of $\mathcal{P}^*$ can improve on one objective without degrading another. NSGA-II
approximates this front by evolving a population of candidate routing plans through
selection, crossover, and mutation~\cite{jafari2021}.

\subsection{DAG Task Model}

Modern IoT applications are not monolithic: they decompose into sequences of
dependent processing steps~\cite{li2025microservice}. The DAG model captures
this as $G = (\mathcal{V}, \mathcal{E})$, where each node $v_i \in \mathcal{V}$
is a microservice step with demand $\langle \ell_i, d_i^{\mathrm{in}},
d_i^{\mathrm{out}} \rangle$ and each directed edge $(v_i, v_j) \in \mathcal{E}$
means $v_j$ cannot start until $v_i$ finishes~\cite{zhu2025dag}. End-to-end
latency along the critical path is:
\begin{equation}
  L_\mathrm{app} = \sum_{v_i \in \mathrm{critical\;path}}
  \bigl[T_\mathrm{exec}(v_i) + T_\mathrm{transfer}(v_i, v_{i+1})\bigr]
  \label{eq:dag}
\end{equation}
where $T_\mathrm{transfer} = 0$ when consecutive steps share the same node.

% =============================================================================
\section{Research Evolution: Four Generations}
\label{sec:evolution}
% =============================================================================

\subsection{Generation 1 — Single-Objective Heuristics (before 2021)}

The earliest task offloading work treated the decision as binary: run locally or
send to the cloud. Greedy algorithms, game theory, and single-objective
metaheuristics such as Particle Swarm Optimization were common~\cite{jafari2021}.
These methods optimized one metric at a time, operated on static topologies, and
assumed devices stayed connected. The core limitation was that energy and latency
pull in opposite directions. Routing everything to the cloud minimizes computation
latency but maximizes transmission energy; running everything locally saves energy
but misses deadlines. A useful solution requires balancing both simultaneously.

\subsection{Generation 2 --- Multi-Objective Evolutionary Algorithms (2021--2023)}

Jafari and Rezvani~\cite{jafari2021} moved the field forward by framing fog
offloading as a multi-objective problem and applying NSGA-II with Minimax
Differential Evolution (MMDE) mutation. Instead of a single optimum, their
approach produces a Pareto front of solutions that trade off energy against latency,
giving operators the choice of operating point. The MMDE mutation operator:
\begin{equation}
  \mathbf{V} = \mathbf{r}_1 + F \cdot (\mathbf{r}_2 - \mathbf{r}_3),\quad
  F = 0.5,\;\mathrm{CR} = 0.9
  \label{eq:mmde}
\end{equation}
uses the difference between two population members to steer mutation toward
productive regions rather than sampling blindly. The result converges faster and
produces more robust solutions than standard polynomial mutation.

The remaining weakness is that the algorithm treats all tasks equally. Heavy
tasks that would never realistically run on a fog node still consume optimization
cycles, slowing convergence and occasionally crashing fog queues when the solution
assigns them there.

\subsection{Generation 3 --- Pre-Filtering and Hybrid Systems (2024)}

Nouri et al.~\cite{nouri2024tof} introduced the TOF-Broker, a classification
layer that runs before NSGA-II. It computes EC~\eqref{eq:ec} for each task and
immediately routes boulders to the cloud, shrinking the search space by 15--30\%.
On a standard IoT workload this yields a 12.18\% energy reduction and only 0.38\%
additional latency overhead. Concurrently, the microservices
community~\cite{li2025microservice} demonstrated that modeling applications as
DAGs enables partial offloading: lightweight steps stay on the fog while heavy
steps go to the cloud, something monolithic task models cannot express.

The remaining problem is rigidity. Both the TOF-Broker and NSGA-II produce static
routing decisions. When fog load shifts during a simulation run, or when a vehicle
moves into a new zone mid-task, neither system can respond.

\subsection{Generation 4 --- Adaptive AI-Driven Systems (2024--2025)}

The current generation replaces static decision-making with learned agents that
update continuously. In vehicular fog settings, DRL-based
offloading~\cite{vfc_drl2024} showed that DQN agents outperform greedy policies
under high load by spreading tasks more intelligently across fog nodes. CSFD~\cite{csfd2025}
embedded a DRL agent inside the SDN controller to predict traffic bursts and
install routing rules before they arrive. MAVTO~\cite{mavto2024} used vehicle GPS
data to pre-position tasks at the next fog zone before the vehicle crossed the
coverage boundary. Zhu et al.~\cite{zhu2025dag} combined all three ideas at the
component level: DAG modeling, NSGA-II offline pre-training, and online meta-RL.

None of these systems addresses all five challenges at once, which is the gap
this paper focuses on.

% =============================================================================
\section{State of the Art}
\label{sec:sota}
% =============================================================================

\subsection{Multi-Objective Optimization: NSGA-II with MMDE}

Jafari and Rezvani~\cite{jafari2021} remain the standard reference for
multi-objective metaheuristic offloading. Their setup uses binary chromosomes
encoding fog-or-cloud decisions per task, evaluated on the energy and latency
objectives defined in equations~(2) and~(3). Non-dominated sorting produces Pareto
fronts across generations; crowding distance maintains diversity within each front;
and MMDE~\eqref{eq:mmde} replaces random mutations with directional ones drawn
from the population's own spread. The practical consequence is that the algorithm
discovers the energy--latency trade-off surface reliably and can be tuned to favor
either objective by selecting a different point on the Pareto front.

What this work does not provide is any mechanism to handle tasks that are
computationally infeasible for fog nodes. All tasks enter the optimization
pipeline regardless of their cost, and the resulting solutions occasionally
assign boulders to fog nodes, causing queue starvation for lighter tasks.

\subsection{Task Pre-Filtering: TOF-Broker}

Nouri et al.~\cite{nouri2024tof} address this directly. The TOF-Broker computes
EC~\eqref{eq:ec} for each incoming task before any optimizer runs and immediately
separates boulders from pebbles. Boulders go to the cloud without entering the
search space; only pebbles are handed to NSGA-II. This produces three concrete
benefits: fog queue crashes disappear, NSGA-II converges faster on a smaller
problem, and overall energy consumption drops by 12\%. The 0.38\% latency
overhead is negligible in practice.

The open issue the TOF-Broker creates is what the authors call the pebble avalanche.
When task arrival rates exceed roughly 1,000 per second, pebbles accumulate in
the queue faster than NSGA-II can optimize them. The broker was designed to prevent
fog crashes, but under extreme load it simply defers the crash rather than
eliminating it.

\subsection{Software-Defined Networking for Fog Traffic}

Gasmi and Ezzedine~\cite{gasmi2024sdn} surveyed SDN-based approaches to fog
offloading and established the theoretical case for decoupling routing intelligence
from physical forwarding. In an SDN architecture, forwarding switches hold no
routing state of their own; they execute rules installed by a centralized
controller via the OpenFlow protocol. The controller sees the entire network at
once and can reprogram all switches simultaneously with a single API call. For
task offloading this matters because it enables VIP lanes: when a batch of pebbles
needs to reach the cloud urgently, the controller can open a dedicated high-priority
path across every intermediate switch in under 5\,ms.

The limitation identified in the same survey is the controller bottleneck. When
many devices request new flow rules simultaneously, the controller becomes
overloaded, adding 8--15\,ms latency per flow. This negates the advantage of SDN
for latency-sensitive workloads and is the problem that motivates embedding an RL
agent inside the controller.

\subsection{Microservice Architecture and DAG Modeling}

Li et al.~\cite{li2025microservice} produced the most comprehensive recent survey
of microservice-oriented resource management in fog and edge computing, covering
136 studies from 2020 to 2024. They organize the field around five pillars: service
placement, resource provisioning, task scheduling and offloading, resource
allocation, and replica selection. Their main finding is that while each pillar
is reasonably well-studied individually, the interactions between them are almost
entirely unexplored. Jointly optimizing across all five simultaneously, which is
what a real system must do, has no published solution.

The DAG formalization used in this paper comes from Zhu et al.~\cite{zhu2025dag},
who define application steps as nodes with per-step computational demands and data
sizes, and inter-step dependencies as directed edges. This formulation enables
partial offloading at step granularity and is the basis for the latency model
in equation~\eqref{eq:dag}.

\subsection{Deep Reinforcement Learning for Task Offloading}

\subsubsection{DRL in Vehicular Fog Networks}

An anonymous study~\cite{vfc_drl2024} is the closest published work to the
compute-placement agent in our proposed architecture. It trains a DQN on a
grid-city simulation where client vehicles offload tasks to nearby parked service
vehicles. The agent observes queue depths, distances, and channel quality at all
candidate nodes and selects a destination via epsilon-greedy policy with experience
replay. The main finding is that DRL distributes load more evenly than greedy
policies and avoids the congestion cascades that greedy methods suffer under high
arrival rates.

Three limitations motivate our extensions: the system optimizes delay only with no
energy term; it has no offline pre-training and therefore suffers cold-start
degradation; and it assumes vehicles remain connected throughout task execution,
ignoring handoff.

\subsubsection{SDN-Integrated DRL: CSFD}

The CSFD framework~\cite{csfd2025} embeds a DRL agent directly inside the SDN
control plane. Rather than waiting for a packet to arrive and then consulting the
controller for a routing rule (8--15\,ms per flow), the agent predicts upcoming
traffic bursts and pre-installs OpenFlow rules before the first packet appears.
The result in a vehicular ad-hoc network evaluation is higher packet delivery
ratio, lower latency, and lower controller CPU load compared to both reactive SDN
and a DRL-only baseline without SDN.

CSFD is the direct prior work for the network-routing agent in our proposed
architecture. Its limitation is that it addresses routing only. Deciding which
fog node should compute a given task is handled by a separate, static mechanism
that the paper does not describe or optimize.

\subsubsection{DAG Offloading with Meta-RL and Offline Pre-training}

Zhu et al.~\cite{zhu2025dag} is the architecturally closest published system to
our proposal. They model each application as a DAG, run NSGA-II offline on
historical batches to generate Pareto-optimal routing plans, convert the plans
into labeled (state, action) pairs, and pre-train a Meta-RL agent via behavioral
cloning. The meta-policy then adapts to new network environments quickly because
it already has a sensible prior rather than starting from random weights. Their
results show consistent improvement over both standalone DRL and standalone NSGA-II
in dynamic network conditions.

The gaps in this work are the absence of mobility modeling and an SDN layer. The
paper's future work section explicitly names both as needed extensions.

\subsection{Mobility-Aware Offloading}

\subsubsection{Reactive Handoff: Dual-Buffer System}

The collaborative offloading framework in~\cite{paper5_mobility} was the first to
explicitly address mid-computation vehicle handoffs. When a vehicle disconnects
from a fog node while a task is running, the framework moves the task from the
normal task queue (NTB) to a high-priority handoff queue (HTB). The HTB preempts
all other work and runs at maximum CPU priority until the task completes. The result
is then held in memory, keyed to the vehicle's identifier, and pushed to the vehicle
the moment it reconnects anywhere in the fog network. This eliminates re-submission
and reduces end-to-end latency by over 50\% in high-mobility scenarios compared to
simply dropping and retransmitting.

The approach is entirely reactive: the system takes no action until the physical
disconnection has already occurred. The time between when a disconnection becomes
predictable and when it actually happens is wasted.

\subsubsection{Proactive Handoff: MAVTO}

MAVTO~\cite{mavto2024} introduced the first proactive mobility framework with
a peer-reviewed evaluation on real vehicular datasets. Given a vehicle's GPS
position, speed, and heading from its spatial telemetry, the system calculates
when the vehicle will exit the current fog zone:
\begin{equation}
  T_\mathrm{exit} = \frac{R_\mathrm{fog} - \mathrm{dist}(v_\mathrm{pos},\,
  f_\mathrm{pos})}{v_\mathrm{closing}}
  \label{eq:texit}
\end{equation}
where $v_\mathrm{closing}$ is the component of vehicle velocity directed toward
the zone boundary. This value is compared against the estimated remaining execution
time to select one of three operating modes: \emph{direct} (task finishes before
exit, do nothing), \emph{predictive} (pre-spin a container on the destination fog
node before the vehicle arrives), or \emph{hybrid} (split computation between
local and edge). Compared to non-mobility-aware baselines, MAVTO improves task
completion success rate by more than 25\%.

The mode selection is rule-based, not learned, and the system has no energy
objective and no SDN routing for the container migration path.

\subsubsection{Spatiotemporal Prediction: Digital Twin with STGNN}

The most recent work in this space~\cite{zhu2025spatiotemporal} combines a Digital
Twin (a continuously synchronized virtual replica of the network), a
Spatio-Temporal Graph Neural Network (STGNN) for demand forecasting, and A3C
multi-agent DRL for task caching. The STGNN processes the vehicular network as a
spatial graph with an LSTM temporal layer, predicting task demand per fog zone
5--10 seconds ahead and enabling proactive resource reservation. This is the
current ceiling for trajectory prediction quality but lacks multi-objective
optimization, task pre-filtering, and energy awareness.

\subsection{Comparative Summary}

Table~\ref{tab:comparison} places all ten surveyed works across six capability
dimensions. The proposed Predictive Cloud-Native Mobile Edge (PCNME) architecture
is the only entry that checks all six.

\begin{table}[H]
\centering
\caption{Surveyed works across six capability dimensions.
  \checkmark\,=\,addressed,\;$\circ$\,=\,partial,\;$\times$\,=\,not addressed.}
\label{tab:comparison}
\renewcommand{\arraystretch}{1.3}
\begin{tabularx}{\textwidth}{@{}L{3.1cm}C{1.1cm}C{1.2cm}C{0.9cm}C{0.9cm}
                              C{0.9cm}C{1.2cm}C{1.0cm}@{}}
\toprule
\textbf{Work} & \textbf{Multi-obj.} & \textbf{Pre-filter}
  & \textbf{DAG} & \textbf{DRL} & \textbf{SDN}
  & \textbf{Mobility} & \textbf{Year} \\
\midrule
Jafari \& Rezvani~\cite{jafari2021}
  & \checkmark & $\times$ & $\times$ & $\times$ & $\times$ & $\times$ & 2021 \\
Nouri et al.~\cite{nouri2024tof}
  & \checkmark & \checkmark & $\times$ & $\times$ & $\times$ & $\times$ & 2024 \\
Gasmi \& Ezzedine~\cite{gasmi2024sdn}
  & $\times$ & $\times$ & $\times$ & $\times$ & \checkmark & $\times$ & 2023 \\
Li et al.~\cite{li2025microservice}
  & $\circ$ & $\times$ & \checkmark & $\circ$ & $\times$ & $\times$ & 2025 \\
Zhu et al.~\cite{paper5_mobility}
  & $\times$ & $\times$ & $\times$ & $\times$ & $\times$ & $\circ$ & 2023 \\
Toopchinezhad \& Ahmadi~\cite{vfc_drl2024}
  & $\times$ & $\times$ & $\times$ & \checkmark & $\times$ & $\times$ & 2024 \\
Rahmatian et al.~\cite{csfd2025}
  & $\times$ & $\times$ & $\times$ & \checkmark & \checkmark & $\times$ & 2025 \\
Zhu et al.~\cite{zhu2025dag}
  & \checkmark & $\times$ & \checkmark & \checkmark & $\times$ & $\times$ & 2025 \\
Zaidi et al.~\cite{mavto2024}
  & $\times$ & $\times$ & $\times$ & $\times$ & $\times$ & \checkmark & 2024 \\
Zhu et al.~\cite{zhu2025spatiotemporal}
  & $\times$ & $\times$ & $\circ$ & \checkmark & $\times$ & \checkmark & 2025 \\
\midrule
\textbf{Proposed PCNME}
  & \checkmark & \checkmark & \checkmark & \checkmark
  & \checkmark & \checkmark & 2026 \\
\bottomrule
\end{tabularx}
\end{table}

% =============================================================================
\section{Open Research Gaps}
\label{sec:gaps}
% =============================================================================

The ten papers surveyed above converge on five gaps that none of them closes
individually, and that no joint system in the literature has closed together. The
first two are the primary contributions of this analysis.

\medskip\noindent\textbf{Gap 1 — The Integration Gap.}
The most striking pattern in Table~\ref{tab:comparison} is that every system
addresses at most three of the six capability dimensions. This is not a coincidence
of publication timing: it reflects a deeper structural difficulty. Compute
placement decisions (which fog node handles a DAG step) and network routing
decisions (which physical path carries that step's data) are interdependent in a
way that current architectures do not account for. A placement that minimizes
compute latency may route a task to a fog node whose upstream link is congested,
adding more delay than it saves. Handling this correctly requires either a joint
optimization over placement and routing simultaneously, which creates an
intractable action space, or an architectural design where two specialized agents
share the network-state information they each need. Neither approach has been
demonstrated at the scale of DAG workloads with vehicle mobility.

\medskip\noindent\textbf{Gap 2 — The Cold-Start Gap.}
Every DRL system surveyed initializes from random weights. This means the agent
spends its early operational life making near-random decisions, a period that spans
thousands of interaction steps in typical training curves. For a vehicular safety
system, near-random placement decisions during cold-start represent an operational
hazard, not merely a performance cost. Zhu et al.~\cite{zhu2025dag} demonstrate
that seeding DRL weights from NSGA-II Pareto solutions via behavioral cloning
eliminates much of this degradation, but only in a setting without vehicle
mobility and without SDN. Whether this pre-training advantage persists across the
full complexity of a city-scale vehicular network, with both mobility and network
management in play, remains an open question.

\medskip\noindent\textbf{Gap 3 — The DAG-Mobility Gap.}
The mobility-aware systems reviewed in Section~4.6 treat tasks as indivisible
units. When a task is instead a DAG with multiple concurrent steps running on
different fog nodes, a vehicle handoff may invalidate some steps but not others
depending on which fog node each step is assigned to. The T$_\text{exit}$
comparison in MAVTO~\cite{mavto2024} is done once per task; it would need to be
done once per DAG step, with the additional complexity that the relevant fog node
differs per step. No published work addresses this.

\medskip\noindent\textbf{Gap 4 — The Pebble Avalanche Gap.}
The TOF-Broker solves the boulder problem but exposes a new failure mode. When
tasks arrive faster than roughly 1,000 per second, the pebble queue grows faster
than NSGA-II can drain it. The queue eventually overflows, causing exactly the
kind of cascade failure the broker was designed to prevent. A mechanism for
graceful degradation in this regime, whether by dynamically lowering the EC
threshold to reclassify heavy pebbles as boulders, or by batching overflow pebbles
into super-tasks routed directly to cloud, has not been proposed or validated.

\medskip\noindent\textbf{Gap 5 — The Trajectory Prediction Gap.}
MAVTO predicts vehicle positions using linear dead-reckoning, which works
reasonably well on straight highway segments but degrades at intersections, during
braking events, or when a vehicle takes an unexpected turn. Zhu et al.~\cite{zhu2025spatiotemporal}
demonstrate that a Spatio-Temporal Graph Neural Network trained on historical
mobility data is substantially more accurate, but they apply it only to task
caching, not to proactive handoff. A system that uses STGNN-quality trajectory
prediction as the input to a proactive container migration mechanism, combined with
Pareto-optimal placement, has not been built.

% =============================================================================
\section{Proposed Architecture}
\label{sec:proposed}
% =============================================================================

The gaps identified above point toward a specific architecture. We call it the
Predictive Cloud-Native Mobile Edge (PCNME) and describe its five components
here. The intent is not to present a finished system but to show that the five
gaps have a coherent joint solution.

\subsection{Component Overview}

\textbf{TOF-Broker and Aggregator} address Gaps~1 and~4. The broker applies
EC classification~\eqref{eq:ec} per DAG step before any optimizer runs. When
the pebble queue exceeds a threshold, the aggregator bundles queued pebbles into
a super-task and routes it to the cloud via an SDN VIP lane, preventing avalanche
without losing the data.

\textbf{TOF-MMDE-NSGA-II Offline Optimizer} addresses Gap~2. The combination of
TOF pre-filtering and MMDE-enhanced NSGA-II~\eqref{eq:mmde} generates Pareto-optimal
routing plans from historical task logs. These plans are decomposed into
(state,\;action) pairs for behavioral cloning of both DRL agents before the
simulation begins.

\textbf{DRL Agent~1} (task placement) addresses Gap~1. A Deep Q-Network
pre-trained on NSGA-II Pareto data and fine-tuned online via experience replay
selects which fog node handles each pebble DAG step. Its 13-dimensional state
vector includes fog CPU loads, queue depths, task EC, bandwidth utilization,
vehicle speed and heading, $T_\mathrm{exit}$, remaining deadline, and cloud
backlog. Actions are the four available fog nodes plus cloud.

\textbf{DRL Agent~2} (SDN routing) addresses Gaps~1 and~4. A second DQN embedded
in the SDN controller, pre-trained on historical city traffic patterns, predicts
upcoming traffic bursts and pre-installs OpenFlow flow rules before the traffic
arrives. It shares network-state telemetry with Agent~1, providing the coupling
that resolves the interdependence identified in Gap~1.

\textbf{Trajectory Predictor and NTB/HTB Buffers} address Gaps~3 and~5. The
predictor evaluates $T_\mathrm{exit}$~\eqref{eq:texit} separately for each DAG
step rather than once per task, applying step-level proactive migration when
$T_\mathrm{exec} > T_\mathrm{exit}$. The NTB/HTB dual-buffer system from~\cite{paper5_mobility}
provides deterministic fallback when trajectory prediction fails.

\subsection{Offline-to-Online Pipeline}

The system operates in two phases. In the offline phase, historical task logs
and mobility traces feed TOF-MMDE-NSGA-II. The resulting Pareto solutions are
converted to labeled examples and used to initialize both agent networks via
supervised learning. In the online phase, Agent~1 makes placement decisions in
under 1\,ms per pebble step; Agent~2 manages SDN rules proactively. Both update
continuously from rewards of the form:
\begin{equation}
  R = -0.5\,L_\mathrm{norm} - 0.3\,E_\mathrm{norm}
    - 0.2\,\mathbf{1}[L > D]\cdot\lambda
\end{equation}
where $\lambda = 10$ for safety-critical steps and $1$ otherwise. The
offline-to-online knowledge transfer is what eliminates cold-start and is the
central novelty of the design.

% =============================================================================
\section{Comparative Discussion}
\label{sec:discussion}
% =============================================================================

\textbf{Against Jafari and Rezvani~\cite{jafari2021}.} The proposed system
retains NSGA-II with MMDE as its offline optimizer, so it subsumes rather than
replaces this work. The TOF-Broker reduces the problem NSGA-II has to solve; the
DRL agents handle decisions NSGA-II is too slow to make in real time; and the
mobility layer addresses vehicular dynamics that the original static setup cannot
model.

\textbf{Against Nouri et al.~\cite{nouri2024tof}.} The TOF-Broker is incorporated
unchanged. What the proposed system adds is MMDE mutation (faster convergence),
two online DRL agents (real-time adaptation), the Aggregator (avalanche prevention),
and GPS-driven mobility management. All four additions directly address limitations
that Nouri et al. identify as future work.

\textbf{Against Zhu et al.~\cite{zhu2025dag}.} This is the closest prior work
structurally. The main differences are the TOF pre-filter applied before NSGA-II
runs, the SDN routing agent (Agent~2), and step-level proactive mobility. The
future work section of~\cite{zhu2025dag} explicitly requests mobility support and
SDN integration.

\textbf{Against CSFD~\cite{csfd2025}.} CSFD's DRL-in-SDN design is adopted
directly as Agent~2. What CSFD does not provide is any compute placement
intelligence. Agent~1 fills that role, and the shared telemetry between the two
agents addresses the coupling Gap~1 identifies.

\textbf{Against MAVTO~\cite{mavto2024}.} The $T_\mathrm{exit}$ formula and
three-mode framework are adopted from MAVTO. The proposed system extends this with
step-level assessment, RL-learned mode selection replacing the rule-based threshold
comparison, and SDN-managed migration paths. The NTB/HTB fallback from~\cite{paper5_mobility}
handles cases where the trajectory prediction fails.

% =============================================================================
\section{Conclusion}
\label{sec:conclusion}
% =============================================================================

This paper reviewed ten works on task offloading in IoT--Fog--Cloud systems,
tracing a clear development from single-objective greedy methods through
multi-objective evolutionary algorithms, hybrid pre-filtering systems, and the
current generation of adaptive AI-driven architectures. The review reveals that
the field has reached a point where each individual challenge, the boulder problem,
cold-start, static optimization, network congestion, and mobility, has a credible
solution in isolation. What does not yet exist is a system that handles all of
them together.

The two primary gaps, the Integration Gap and the Cold-Start Gap, define the
most important open problem: whether a dual-agent architecture combining offline
Pareto pre-training with online reinforcement learning can operate reliably at
city scale with DAG workloads and vehicle mobility simultaneously in play. The
evidence from the surveyed works suggests this is achievable by composition; what
is needed is the engineering effort to combine the components and a rigorous
empirical validation.

The proposed PCNME architecture is a concrete step in that direction. Its novelty
is not in any individual component, each of which already exists in the literature,
but in their integration and in the offline-to-online knowledge transfer mechanism
that eliminates cold-start degradation before the system begins operating.

Future directions include replacing linear dead-reckoning with STGNN-based
trajectory prediction~\cite{zhu2025spatiotemporal}, applying federated learning
to train Agent~2 across multiple cities without sharing raw traffic data, and
investigating TD3 or SAC as higher-performing replacements for DQN in the
continuous-action resource allocation setting.

% =============================================================================
\begin{thebibliography}{11}

\bibitem{jafari2021}
V.~Jafari and M.\,H.~Rezvani,
``Joint optimization of energy consumption and time delay in {IoT}--fog--cloud
  computing environments using {NSGA-II} metaheuristic algorithm,''
\textit{Journal of Ambient Intelligence and Humanized Computing},
vol.~14, pp.~1675--1698, 2023 (published online 2021).
\newblock \url{https://link.springer.com/article/10.1007/s12652-021-03388-2}

\bibitem{nouri2024tof}
N.~Nouri and M.\,H.~Rezvani,
``A multi-objective approach for optimizing {IoT} applications offloading in
  fog--cloud environments with {NSGA-II},''
\textit{The Journal of Supercomputing}, 2024.
\newblock \url{https://link.springer.com/article/10.1007/s11227-024-06431-z}

\bibitem{gasmi2024sdn}
K.~Gasmi and T.~Ezzedine,
``Reinforcement learning based task offloading of {IoT} applications in fog
  computing: algorithms and optimization techniques,''
\textit{Cluster Computing}, 2024.
\newblock \url{https://www.ijcnis.org/index.php/ijcnis/article/view/8458/2522}

\bibitem{li2025microservice}
Y.~Li et al.,
``Energy-efficient resource management in microservices-based fog and edge
  computing: state-of-the-art and future directions,''
\textit{ACM Computing Surveys}, 2025.
\newblock \url{https://arxiv.org/pdf/2512.04093}

\bibitem{paper5_mobility}
M.~Zhu, J.~Liu, and H.~Wang,
``A collaborative offloading task framework for {IoT} fog computing,''
in \textit{Proc.\ International Conference on Fog and Mobile Edge Computing},
2023.

\bibitem{vfc_drl2024}
M.\,P.~Toopchinezhad and M.~Ahmadi,
``Deep reinforcement learning for delay-optimized task offloading in vehicular
  fog computing,''
\textit{arXiv preprint arXiv:2410.03472}, 2024.
\newblock \url{https://arxiv.org/html/2410.03472}

\bibitem{csfd2025}
M.~Rahmatian, A.~Javadpour, and F.~Ja'fari,
``{CSFD}: Enhancing congestion reduction in vehicular communication with {SDN}
  and fog integration using deep reinforcement learning,''
\textit{Analog Integrated Circuits and Signal Processing}, 2025.
\newblock \url{https://link.springer.com/article/10.1007/s10470-025-02541-7}

\bibitem{zhu2025dag}
L.~Zhu, T.~Liu, Y.~Yang, and X.~Luo,
``Dynamic task offloading scheme for edge computing via meta-reinforcement
  learning,''
\textit{Journal of Manufacturing Systems}, 2025.
\newblock \url{https://www.sciencedirect.com/science/article/pii/S1546221825000712}

\bibitem{mavto2024}
S.~Zaidi, M.\,A.~Attalah, L.~Khamer, and C.\,T.~Calafate,
``Task offloading optimization using {PSO} in fog computing for the internet
  of drones,''
\textit{Drones}, vol.~8, no.~11, p.~696, 2024.
\newblock \url{https://www.mdpi.com/2504-446X/8/11/696}

\bibitem{zhu2025spatiotemporal}
L.~Zhu et al.,
``A vehicular edge computing offloading and task caching solution based on
  spatiotemporal prediction,''
\textit{Future Generation Computer Systems}, vol.~166, p.~107679, 2025.
\newblock \url{https://www.sciencedirect.com/article/abs/pii/S0167739X24006435}

\bibitem{mashal2024drl}
I.~Mashal et al.,
``Multiobjective offloading optimization in fog computing using deep
  reinforcement learning,''
\textit{Journal of Computer Networks and Communications}, 2024.

\end{thebibliography}

\end{document}
